\sscp{\tsc{East Sumbawa} and Komodo}{03-01-01}
As mentioned above, the main changes which are shared between
Bima, Kolo, and Sanggar also occur in Komodo,
as well as some other \tsc{Bima-Lembata} languages.
These changes are *R > {\0}, * ay > \it{e,}
and *-aw > \it{o}.
In addition to the links between \tsc{East Sumbawa}
and Komodo, there are also several changes exclusively shared
between Bima and \tsc{Havu-Dhao} (\srf{03-03-03}).
The multiple links between Bima
and other languages may be due to contact.
Bima had considerable influence in the wider region,
particularly after the influence of the Sultan of Goa (Makassar)
waned from 1727 onwards.

The change *R > {\0} is shared by \tsc{East Sumbawa} and Komodo.
Although \tsc{Sumba-Havu} often also has *R > {\0},
there is evidence that in \tsc{Sumba-Havu} *R merged with *y
after which it was lost.
That is \tsc{Sumba-Havu} usually has *R/*y > **y > {\0} (see \srf{03-03},
particularly \trf{03-51} on \prf{tab:03-51}).\footnote{
		There are also a handful of forms with *R = \it{r} in \tsc{Sumba}
		where it is lost in \tsc{East Sumbawa}; e.g.
		*paRih ‘stingray’, *uRat ‘vein’, and *baqəRu ‘new’ in \trf{03-09}.}
In Bima *y > \it{ʤ}
and is thus distinct compared with *R > {\0}.
Examples of *y > \it{ʤ} in Bima include: *bayu > \it{mbaʤu} ‘pound’,
*layaR > \it{loʤa} ‘sail’,
*duyuŋ > \it{ruʤu} ‘dugong’,
*kahiw > **kayu > \it{haʤu} ‘tree,
wood’, and *ma-həyaq > \it{maʤa} ‘ashamed,
shy’ \citep[80]{Blust2008}.\footnote{
		Havu has one possible instance
		of PMP *y > \it{ʤ}; *duyuŋ > \it{ruʤu} ‘dugong’ \citep[332]{Jonker1908}.
		In other cases it has *y > {\0},
		often with fronting of an adjacent vowel.
		Two examples include *layaR > \it{lai} ‘sail’
		and *buqaya › \it{voe} ‘crocodile’.
		\tsc{Havu-Dhao} also has secondary **y > \it{ʄ}, \it{ʤ}, \it{y}.
		Examples include *kahiw > **kayu > Havu,
		Dhao \it{aʄu} ‘wood’ and *i-aku > Havu (Dimu,
		Raijua Wawa) \it{ʄaa}, Havu (Raijua D{\Q}ida) \it{ʤoo},
		Havu (Seba, Liae, Mehara) \it{yaa}, Dhao \it{ʤaʔa} \tsc{‘1sg’}.}
Thus, the changes affecting PMP *R in \tsc{East Sumbawa}
and \tsc{Sumba-Havu} cannot be taken as shared.

\begin{table}[p]\centering\tcp{\tsc{East Sumbawa} and Komodo *R > {\0}}{03-09}
\begin{adjustbox}{width=\textwidth}
	\begin{threeparttable}\st{2pt}
	\begin{tabular}{lllllllll}\lsptoprule
*gloss	&	house	&	root	&	run	&	red	&	thorn\su{iii}	&	rotten	&	stingray	&	vein\\
PMP	&	*\tbr{R}umaq	&	*\tbr{R}amut	&	*la\tbr{R}iw	&	*ma-i\tbr{R}aq	&	*du\tbr{R}i	&	*ba\tbr{R}iw	&	*pa\tbr{R}ih\su{iv}	&	*u\tbr{R}at\\  \midrule
Sanggar	&	{\hr}uma	&	\cg{(aka)}	&	{\hr}mpalai	&	{\hr}mea	&	{\hr}rui	&	{\hr}wai	&	&{\hr}ua\\
Kolo	&	{\hr}uma	&	&	&	&	&	&	&	\\
Bima	&	{\hr}uma	&	{\hr}amu	&	{\hr}rai	&	\cg{(kala)}	&	{\hr}rui	&	{\hr}mbai	&	{\hr}fai	&	{\hr}kaʔuʔa\\
Komodo	&	\cg{(dekaŋ)}	&	\cg{(wakeʔ)}	&	{\hr}pelai	&	\cg{(rara)}	&	\cg{(toko)}	&	\cg{(wueh)}	&	{\hr}pai	&	{\hr}uraʔ{\cgr}\\  \midrule
Manggarai	&	\cg{(mabru)}	&	{\hr}ramut\su{i}	&	\cg{(ləmpok)}	&	\cg{(ʋara)}	&	\cg{(toko)}	&	&{\hr}pai	&	{\hr}urak\\
Lio	&	\cg{(saʔo)}	&	{\hr}kamu\su{ii}	&	{\hr}paru	&	{\hr}mera	&	\cg{(toko)}	&	\cg{(məʋu)}	&	{\hr}pari	&	{\hr}ura\\
Sika	&	\cg{(ləpo)}	&	{\hr}ʔramut	&	{\hr}plari	&	{\hr}merak	&	{\hr}luriŋ	&	\cg{(bewus)}	&	{\hr}paʔe, pəri	&	\cg{(aliŋ)}\\
Kambera	&	{\hr}uma	&	{\hr}amu	&	{\hr}palai	&	{\hr}mii	&	{\hr}rii	&	\cg{(muhuŋ)}	&	{\hr}pai, pari	&	{\hr}ura, urat u\\
Havu	&	{\hr}əmu	&	{\hr}amo	&	{\hr}perai	&	{\hr}mea	&	{\hr}rui	&	&	&	\cg{(lua)}\\
Dhao	&	{\hr}əmu	&	{\hr}amo	&	{\hr}rai	&	{\hr}mea	&	{\hr}rui	&	&{\hr}ai	&	\cg{(kalua)}\\
\midrule\midrule
	\end{tabular}
	\begin{tabular}{lllllllll}
*gloss	&	new	&	blood	&	needle	&	post	&	brave	&	bathe	&	give	&	vine\\
PMP	&	*baqə\tbr{R}u	&	*da\tbr{R}aq	&	*za\tbr{R}um	&	*hadi\tbr{R}i	&	*ba\tbr{R}ani	&	*di\tbr{R}us	&	*bə\tbr{R}ay	&	*wa\tbr{R}əj\su{ix}\\ \midrule
Sanggar	&	{\hr}weu	&	‹rah›	&	{\hr}rao	&	&{\hr}wani	&	{\hr}ndiu	&	\<bèkè\>\su{viii}	&	\cg{(ika)}\\
Kolo	&	&	&	&	&	&{\hr}riu	&	{\hr}weya	&	{\hr}au\\
Bima	&	{\hr}ɓou	&	{\hr}raʔa	&	{\hr}nda(ʔ)u\su{v}	&	{\hr}riʔi	&	{\hr}mbani	&	{\hr}ndeu	&	{\hr}mbei	&	{\hr}ai\\
Komodo	&	{\hr}wou	&	{\hr}raa	&	{\hr}rao	&	{\hr}kompo rii	&	{\hr}bani\su{vii}	&	\cg{(buri)}	&	{\hr}wiaŋ	&	\cg{(tali)}\\  \midrule
Manggarai	&	{\hr}ʋəru	&	{\hr}dara	&	{\hr}ʤaruŋ\su{vi}	&	{\hr}siri	&	{\hr}rani	&	\cg{(ʧəboŋ)}	&	\cg{(sotor)}	&	\cg{(ʋase)}\\
Lio	&	\cg{(muri)}	&	{\hr}raa	&	{\hr}laru	&	\cg{(ləke)}	&	{\hr}raniŋ	&	\cg{(huʔi)}	&	{\hr}bəli	&	\\
Sika	&	{\hr}βəruŋ	&	\cg{(mei)}	&	\cg{(luhir)}	&	{\hr}liri	&	\cg{(bərani)}	&	\cg{(həbo)}	&	\cg{(nein)}	&	\cg{(taleʔ)}\\
Kambera	&	{\hr}mbaru	&	{\hr}ria	&	\cg{(utu)}	&	\cg{(punduk u)}	&	{\hr}mbeni	&	\cg{(ihu)}	&	{\hr}wua	&	\cg{(liku)}\\
Havu	&	{\hr}viu	&	{\hr}raa	&	{\hr}ʄau	&	{\hr}gerii	&	{\hr}ɓani	&	{\hr}ʄiu	&	{\hr}vie	&	\cg{(dari)}\\
Dhao	&	{\hr}hiu	&	{\hr}raa	&	{\hr}ʄaʄau	&	{\hr}agarii	&	{\hr}bani	&	{\hr}diu	&	{\hr}hia	&	\cg{(ɖ͡ʐari)}\\
		\lspbottomrule
	\end{tabular}
		\begin{tablenotes}\tnsize
			\item[i]	{Manggarai \it{ramut} is ‘moss’.}
			\item[ii]	{Lio \it{kamu} ‘root’ is from earlier **kramu \citep[116]{Elias2018}.}
			\item[iii]	{Sanggar \it{rui}
									and Kambera \it{rii} are ‘thorn,
									bone’, Bima \it{rui} is ‘thorn’,
									Sika \it{luriŋ}, as well as Havu
									and Dhao \it{rui}
									are ‘bone’.
									Komodo, Manggarai, and Lio \it{toko} are ‘bone’.
									Komodo has \it{dəkeʔ} ‘thorn’,
									Manggarai has \it{karot} ‘thorn’,
									and Lio has \it{gulu} ‘thorn’,
									as well as \it{karo} ‘long sharp thorn.’}
			\item[iv]	{\citet[vol. 2; 348]{Verheijen1995}
								gives Manggarai \it{pai} ‘stingray’ as “probably a loan from Bimanese”.
								Lio \it{pari} is ‘a big ocean fish’ \citep[356]{Arndt1933}.
								Sika \it{paʔe} is from \citet[155]{PareiraLewis1998}
								and \it{pəri} from \citet[vol. 2; 347]{Lewis1995}.
								\it{pəri} does not occur in \citet[155]{PareiraLewis1998}.}
			\item[v]	{Bima \it{ndau} is given as ‘needle’ and \it{ndaʔu} as ‘sew’ \citep[99]{Ismail1985}.}
			\item[vi]	{Manggarai \it{ʤaruŋ} ‘needle’ is probably a loan given irregular *z > \it{ʤ}
								(cf. Malay \it{ʤarum}).
								Rembong \it{zarun} ‘needle’ with regular *z > \it{z} shows \tsc{West
								Flores} *R = \it{r.}}
			\item[vii]	{Komodo \it{bani} is ‘angry’.}
			\item[viii]	{Sanggar \it{\<bèkè\>} is one of three words which shows unexpected
								\it{k} between two identical vowels.
								The other two examples are *dahun > \it{rokon} ‘leaf’ (Bima \it{roʔo})
								and *puqun > \it{pukun} ‘tree trunk’ (Bima \it{fuʔu}).}
			\item[ix]	{Putative reflexes of *waRəj here mean ‘rope’.
								Bima also has *talih > \it{dari} ‘hanging rope’.}
		\end{tablenotes}
	\end{threeparttable}
\end{adjustbox}
\end{table}

Examples of *R > {\0} in \tsc{East Sumbawa}
and Komodo are given in \trf{03-09}, along with cognates from each of the other branches
of \tsc{Bima-Lembata} for comparison.
While some words in Bima appear to show *R > \it{ʔ},
this is most likely due to a latter process of intervocalic glottal stop insertion.
This is exemplified more clearly in words such as *niuR > \it{niʔu}
‘coconut’, *dahun > **roo > \it{roʔo} ‘leaf’,
and *qauR > **oo > \it{oʔo} ‘bamboo’.

There are also a handful of instances of *R > \it{r} in \tsc{East Sumbawa} and/or Komodo.
Examples include *baRu > Sanggar,
Komodo \it{waru} ‘sea hibiscus’ (Bima \it{wau}),
*bəRəqat > Bima \it{bara} ‘heavy’ (Sanggar \it{dodo},
Komodo \it{mədeh} [both non-cognate]),
and *təRas > Bima \it{tera} ‘hard,
heartwood’, Kolo ‹katòo› [katoʔo] ‘hard’ (Sanggar ‹ka-ekh›
‘hard’ appears to be non-cognate).

The changes affecting final vowel-glide sequences *-ay
and *-aw are shared between \tsc{East Sumbawa},
Komodo, \tsc{Havu-Dhao}, and -- in the case of *-aw > \it{o} -- \tsc{West Flores}.
Sika and Kedang also share *-ay > \it{e}
and *-aw > \it{o}.
But other \tsc{Flores-Lembata} languages have
a variety of reflexes which show that the full vowel-glide sequence
must be reconstructed for Proto-Bima-Lembata.
\tsc{Central Flores} also has sporadic examples of *-ay > \it{e}
and *-aw > \it{o}.
Examples of *-ay > \it{e}
and *-aw > \it{o} are given in \trf{03-11}
and \trf{03-12} respectively.

\begin{table}\tcp{\tsc{Bima-Lembata} *-ay}{03-11}
\begin{adjustbox}{width=\textwidth}
	\begin{threeparttable}\st{2pt}
	\begin{tabular}{llllllll}\lsptoprule
*gloss&rice plant&liver&die, dead&male&rattan&oar&female\\
PMP&*paj\tbr{ay}&*qat\tbr{ay}&*mat\tbr{ay}&*ma-Ruqan\tbr{ay}&*qu\tbr{ay}&*bəRs\tbr{ay}\su{‡}&*bin\tbr{a}h\tbr{i}\\ \midrule
Sanggar&{\hr}par\tbr{e}&{\hr}at\tbr{e}&{\hr}mat\tbr{e}&{\hr}mon\tbr{e}&{\hr}wu\tbr{e}&\cg{(krawe)}&{\hr}win\tbr{e}\\
Kolo&{\hr}par\tbr{e}&{\hr}at\tbr{e}&{\hr}mat\tbr{e}&&&&\\
Bima&{\hr}far\tbr{e}&{\hr}ad\tbr{e}&{\hr}mad\tbr{e}&{\hr}mon\tbr{e}&&{\hr}wes\tbr{e}&{\hr}win\tbr{e}\\
Komodo&{\hr}par\tbr{e}&{\hr}at\tbr{e}&{\hr}mat\tbr{e}&{\hr}mon\tbr{e}&{\hr}u\tbr{e}&\cg{(kərawe)}&{\hr}win\tbr{e}\\
Havu&{\hr}ar\tbr{e}&{\hr}ad\tbr{e}&{\hr}mad\tbr{e}&{\hr}mon\tbr{e}&&\cg{(tuku)}&{\hr}bən\tbr{i}\\
Dhao&{\hr}ar\tbr{e}&{\hr}aɖ͡ʐ\tbr{e}&{\hr}maɖ͡ʐ\tbr{e}&{\hr}mon\tbr{e}&&{\hr}seh\tbr{e}&{\hr}b͡βən\tbr{i}\\ \midrule
Manggarai&\cg{(ʋoʤa)}&{\hr}at\tbb{i}&{\hr}mat\tbb{a}&{\hr}ron\tbb{a}&{\hr}ʋu\tbb{a}&\cg{(bise)}&{\hr}ʋin\tbb{a}\\
Lio&{\hr}par\tbr{e}\su{†}&{\hr}at\tbr{e}&{\hr}mat\tbb{a}&\cg{(kaki)}&{\hr}u\tbb{a}&{\hr}ʋes\tbb{a}&\cg{(fai)}\\
Ile Mandiri&{\hr}par\tbr{e}&{\hr}at\tbr{e}&{\hr}mat\tbb{a}&\cg{(ama lake)}&{\hr}u\tbb{a}, u\tbb{ar}&{\hr}bah\tbb{a}&\cg{(ina wae)}\\
Kambera&{\hr}par\tbb{i} \it{‘rice’}&{\hr}et\tbb{i}&{\hr}met\tbb{i}&{\hr}min\tbb{i}&{\hr}iw\tbb{i}&{\hr}buh\tbb{i}&{\hr}kawin\tbb{i}\\
		\lspbottomrule
	\end{tabular}
		\begin{tablenotes}\tnsize
			\item[†]	{
								Lio \it{pare} ‘rice plant’ has irregular medial *j > \it{r}
								(expect *j > \it{ʤ})
								and may be a borrowing.}
			\item[‡]	{Dhao \it{sehe} with consonant metathesis
								is probably from a Rote language,
								all of which have the same metathesis with \it{sefe} `oar'.
								Manggarai \it{bise} ‘paddle’ is a
								Makasar loan \citep[vol. 2; 527]{Verheijen1995}.}
			%\item[Δ]	{Havu \it{bəni}
									%and Dhao \it{b͡βəni} ‘woman’ may be regular from **bine as both
									%have metathesis of a high vowel with a following non-high vowel (\srf{03-03}).
									%No other reflexes of words shaped *(C)V\tsc{[+high]}Cay or *(C)V\tsc{[+high]}Caw
									%have so far been attested in \tsc{Havu-Dhao}.}
		\end{tablenotes}
	\end{threeparttable}
\end{adjustbox}
\end{table}

\begin{table}[p]\centering\tcp{\tsc{Bima-Lembata} *-aw}{03-12}
\begin{adjustbox}{width=\textwidth}
	\begin{threeparttable}\st{2pt}
	\begin{tabular}{llllllll}\lsptoprule
*gloss&sun, day&go&mouse, rat&upper&sugar palm&scratch&steal\\
PMP&*qaləj\tbr{aw}&*lak\tbr{aw}&*balab\tbr{aw}&*bab\tbr{aw}&*qanah\tbr{aw}&*kaR\tbr{aw}&*panak\tbr{aw},\\
&&&&&&&*tak\tbr{aw}\\ \midrule
Sanggar&&{\hr}lak\tbr{o}&\cg{(teku)}&&{\hr}na\tbr{o}&&\cg{(maliŋ)}\\
Kolo&\cg{(ando)}&{\hr}la\tbr{o}&&&&&\\
Bima&{\hr}lir\tbr{o}\su{†}&{\hr}la\tbr{o}&{\hr}karaw\tbr{o}&{\hr}waw\tbr{o}&{\hr}na\tbr{o}&{\hr}ka\tbr{o}&\cg{(mpaŋa)}\\
Komodo&{\hr}ro\tbr{o}&{\hr}lah\tbr{o}&{\hr}bəlaw\tbr{o}&{\hr}wew\tbr{o}&{\hr}na\tbr{o}&{\hr}ha\tbr{o}&{\hr}kenah\tbr{o}\\
Havu&{\hr}loɗ\tbr{o}&{\hr}kak\tbr{o}&\cg{(keɓuku)}&\cg{(ɗida)}&&{\hr}ka\tbr{o}, ka\tbr{u}&{\hr}menaʔ\tbr{o}\\
Dhao&{\hr}loɗ\tbr{o}&{\hr}kak\tbr{o}&{\hr}marah\tbr{o}&\cg{(deɖ͡ʐa)}&&{\hr}ka\tbr{o}&{\hr}manaʔ\tbr{u}\\
Manggarai&{\hr}ləs\tbr{o}&{\hr}lak\tbr{o}&{\hr}laʋ\tbr{o}&\cg{(eta)}&\cg{(tuak)}&\cg{(koer)}&{\hr}tak\tbr{o}\\ \midrule
Lio&{\hr}ləʤ\tbb{a}&{\hr}laʔ\tbb{a}&\cg{(teʔu)}&{\hr}waw\tbr{o}&{\hr}naa, na\tbr{o}\su{‡}&\cg{(sasi)}&{\hr}nak\tbb{a}\\
Ile Mandiri&{\hr}rər\tbb{a}&\cg{(pana)}&\cg{(kərome)}&\cg{(lolon)}&\cg{(kebo)}&{\hr}gaʔ\tbb{u}&{\hr}tak\tbb{a}ʔ\\
Kambera&{\hr}loɗ\tbb{u}&{\hr}lak\tbb{u}&{\hr}kala\tbb{u}&&&{\hr}ka\tbb{u}&\cg{(maŋaŋa)}\\
		\lspbottomrule
	\end{tabular}
		\begin{tablenotes}\tnsize
			\item[†]	{Bima \it{liro} ‘sun’ may not be from *qaləjaw.
								It may instead be cognate with another set of words in Bima-Lembata that mean ‘sky’.
								Examples include: Lio \it{liru} ‘sky’ (alongside Lio *qaləjaw
								> \it{ləʤa} ‘sun, day’) and Havu \it{liru} ‘sky’ (alongside Havu *qaləjaw > \it{loɗo} ‘day, sun’).}
			\item[‡]	{\citet{Arndt1933} gives \it{naa} ‘a kind of tree’
								and \it{nao} ‘fibres of the raffia palm,
								rope, rope made from fibres’ either of which could reasonably
								be reflexes of *qanahaw.}
		\end{tablenotes}
	\end{threeparttable}
\end{adjustbox}
\end{table}

Bima and Kolo each have a split of *k > \it{h,
k}, {\0}, but in amongst the (limited) Kolo data there are two
words which have differing reflexes compared with Bima: *kədəŋ
> Kolo \it{horu}, Bima \it{kidi} ‘stand’ and *kaRat > Kolo \it{kaka}, Bima \it{haʔa} ‘bite’.
Bima *k > \it{h,
k} also finds parallels in Komodo.
Whenever Komodo has medial *-k- > \it{h},
Bima also has *k > \it{h}.
However, the reverse does not hold
and there are words with medial *k > \it{h} in Bima,
but *k = \it{k} in Komodo.
Examples of medial *k are given in \trf{03-13}.

\begin{table}[p]\centering\tcp{\tsc{East Sumbawa}, Komodo, and Manggarai *k /V{\gap}V}{03-13}
\begin{adjustbox}{width=\textwidth}
	\begin{threeparttable}\st{2pt}
	\begin{tabular}{lllllllll}\lsptoprule
*gloss&go&male&afraid&\tsc{1sg}&fingernail&elbow&bat&civet\\
PMP&*la\tbb{k}aw&*la\tbb{k}i&*ta\tbb{k}ut&*a\tbb{k}u&*kuh\tbb{k}uh&*si\tbb{k}u&*pani\tbb{k}i&**la\tbb{k}u\su{Δ}\\ \midrule
Sanggar&{\hr}la\tbb{k}o&&\cg{(podo)}&{\hr}a\tbb{k}u&{\hr}ku\tbb{k}u&{\hr}si\tbb{k}u&{\hr}ni\tbb{k}i&\\
Kolo&{\hr}lao{\cgr}&&&‹\tbr{h}aü›{\cgr}&{\hr}u\tbr{h}u{\cgr}&&&\\
Bima&{\hr}lao{\cgr}&{\hr}ra\tbr{h}i{\cgr}&{\hr}da\tbr{h}u{\cgr}&{\hr}na\tbr{h}u{\cgr}&{\hr}u\tbr{h}u{\cgr}&{\hr}ʧi\tbr{h}u, ʧi\tbb{k}u\su{‡} {\cgr}&{\hr}pani\tbr{h}i{\cgr}&{\hr}ra\tbr{h}u{\cgr}\\
Komodo&{\hr}la\tbr{h}o{\cgr}&{\hr}la\tbr{h}i{\cgr}&{\hr}nda\tbr{h}u{\cgr}&{\hr}a\tbr{h}u{\cgr}&{\hr}wu\tbb{k}u&{\hr}si\tbb{k}u&{\hr}ni\tbr{t}i&{\hr}la\tbb{k}u\\
Manggarai&{\hr}la\tbb{k}o&{\hr}la\tbb{k}i&\cg{(kəkop)}&{\hr}a\tbb{k}u&{\hr}ʋu\tbb{k}u&{\hr}ʧi\tbb{k}u&{\hr}ne\tbb{k}i\su{†}&{\hr}\tbb{k}ula\\
		\lspbottomrule
	\end{tabular}
		\begin{tablenotes}\tnsize
			\item[†]	{Manggarai \it{neki} is the general word for ‘several species
								of bat; \it{Pteropus} sp
								and \it{Dobosonia} \it{peroni}’ \citep[375]{Verheijen1967}.
								The Riwu and Sita varieties of Manggarai have \it{niki} ‘bat’.}
			\item[‡]	{\it{ʧihu} from \citet[21]{Ismail1985},
								\it{ʧiku} from \citet[203]{Zollinger1850}.}
			\item[Δ]	{Regional cognates of **laku ‘civet’ include Sika \it{laʔu},
								Kalikasa \it{lako}, Welaun (\tsc{Central Timor}) \it{laʔu},
								and Waima{\Q}a (\tsc{Timor-Babar}) \it{laku}.
								Manggarai \it{kula} ‘civet’ may be cognate with metathesis of
								the final and initial syllables -- such metathesis is used as secret language in Manggarai \citep[460]{Verheijen1941}.
								\citet[245]{Verheijen1967} also lists the word \it{laku} which
								‘still occurs in the name of large thorny plants’
								and connects it with \it{kula} ‘civet’.}
		\end{tablenotes}
	\end{threeparttable}
\end{adjustbox}
\end{table}

Word initially Bima, Komodo,
and Manggarai also have *k > \it{h,
k,} but in this word position the reflexes show even less agreement.
While some cognates sets do share *k- > \it{h-} in all three
languages, there are also cases where one language has *k > \it{h}
while another has *k = \it{k}.
Examples are given in \trf{03-14}.

\begin{table}[t]\centering\tcp{\tsc{East Sumbawa}, Komodo, and Manggarai *k /{\#\gap}}{03-14}
\begin{adjustbox}{width=\textwidth}
	\begin{threeparttable}\st{3pt}
	\begin{tabular}{lllllllll}\lsptoprule
*gloss&lightning&stand&scratch&hook&WB&skin&louse&wood\\
PMP&*\tbb{k}ilat&*\tbb{k}ədəŋ&*\tbb{k}aRaw&*\tbb{k}awit&**\tbb{k}eda\su{†}&*\tbb{k}ulit\su{‡}&*\tbb{k}utu&*\tbb{k}ahiw\\ \midrule
Sanggar&&\cg{(tendi)}&&&&{\hr}lo\tbb{k}i&{\hr}\tbb{k}utu&{\hr}\tbb{k}aʤu\\
Kolo&&{\hr}\tbr{h}oru{\cgr}&&&&{\hr}kalu\tbb{k}i&&{\hr}\tbr{h}ai{\cgr}\\
Bima&{\hr}\tbb{k}ila&{\hr}\tbb{k}idi&{\hr}\tbb{k}ao&{\hr}\tbr{h}awi{\cgr}&{\hr\hr}\tbr{h}era{\cgr}&{\hr}\tbr{h}uri{\cgr}&{\hr}\tbr{h}udu{\cgr}&{\hr}\tbr{h}aʤu{\cgr}\\
Komodo&{\hr}\tbb{k}ila&{\hr}\tbr{h}əre{\cgr}&{\hr}\tbr{h}ao{\cgr}&&{\hr\hr}\tbb{k}esa&{\hr}lu\tbb{k}i&{\hr}\tbr{h}utu{\cgr}&{\hr}\tbr{h}aʤu{\cgr}\\
Manggarai&{\hr}\tbr{h}ilat{\cgr}&{\hr}\tbr{h}əse{\cgr}&\cg{(ŋgetit)}&{\hr}\tbb{k}aet&{\hr\hr}\tbb{k}esa&{\hr}huki&{\hr}\tbr{h}utu&{\hr}\tbr{h}aʤu{\cgr}\\
		\lspbottomrule
	\end{tabular}
		\begin{tablenotes}\tnsize
			\item[†]	{**keda has many regional cognates including Sika \it{kera},
								Kambera \it{yera,} Proto-Rote-Meto **kera \citep{Edwards2021}, and Helong \it{kela}.
								Komodo \it{kesa} is probably a loan from Manggarai,
								as shown by medial **d > \it{s} rather than expected **d > \it{r}.}
			\item[‡]	{Sanggar \it{loki},
								Kolo \it{kaluki} and Komodo \it{luki} have metathesis of the
								initial and medial consonants.
								Such metathesis is also attested in Sumbawa [smw]
								and Waerana with \it{lukit} ‘skin’ which preserves the final consonant.
								If Manggarai \it{huki} ‘skin’ is a reflex of *kulit it has metathesis
								of the medial and final consonants and irregular *t > \it{k}.}
		\end{tablenotes}
	\end{threeparttable}
\end{adjustbox}
\end{table}
